% Requires:
%   \usepackage{tikz}
%   \usetikzlibrary{positioning, arrows.meta, shapes.geometric, fit, backgrounds}
%   \usepackage{fontawesome}     % for \faUser; or replace with text "U"

\begin{figure*}[t]
\centering
\begin{tikzpicture}[
    user/.style={circle, draw=black!70, thick, fill=black!5, minimum size=0.9cm, font=\small},
    privval-strike/.style={rectangle, rounded corners=2pt, draw=black!25, thick,
        minimum width=1.0cm, minimum height=0.45cm, font=\scriptsize, inner sep=2pt,
        text=black!30},
    sanitizer/.style={rectangle, rounded corners=4pt, draw=blue!60!black, very thick,
        fill=blue!8, minimum width=1.4cm, minimum height=0.7cm,
        font=\small\bfseries, text=blue!50!black, align=center},
    cached/.style={circle, draw=blue!70!black, thick, fill=blue!10,
        minimum size=0.7cm, font=\scriptsize},
    derived/.style={rectangle, rounded corners=3pt, draw=gray!60, fill=gray!8,
        font=\scriptsize, inner sep=3pt},
    tgt/.style={rectangle, rounded corners=3pt, draw=orange!70, fill=orange!8,
        font=\scriptsize, inner sep=3pt},
    service/.style={rectangle, rounded corners=4pt, draw=red!60!black, very thick,
        fill=red!5, minimum width=2.4cm, minimum height=1.2cm,
        font=\small\bfseries, text=red!50!black, align=center},
    same/.style={draw=black!30, thick, dashed, rounded corners=2pt},
    arr/.style={-{Stealth[length=4pt]}, thick},
    flow/.style={-{Stealth[length=5pt]}, ultra thick, black!50},
    boundary/.style={dashed, very thick, red!50!black},
    smalllabel/.style={font=\scriptsize, text=black!55},
]

% =====================================================================
%  USER (left): icon + private value bank (struck through)
% =====================================================================

\node[user] (user) at (0, -0.2) {\faUser};

\node[privval-strike] (v1) at (0, -1.45) {$x_1$};
\node[privval-strike] (v2) at (0, -1.95) {$x_2$};
\node[privval-strike] (v3) at (0, -2.45) {$x_3$};

\begin{scope}[on background layer]
    \node[draw=black!40, thick, fill=black!2, rounded corners=3pt,
          fit=(v1)(v3), inner xsep=4pt, inner ysep=4pt] (vbank) {};
\end{scope}
\node[smalllabel, anchor=north] at (vbank.south) {originals deleted};
\node[font=\small\bfseries, text=black!60, anchor=south] at (user.north) {User};

% Initialization arrow: user -> agent (passes values + sanitizer choice)
\draw[arr, black!50] (vbank.east) -- node[smalllabel, above, align=center]
    {initialize:\\$X, \mathcal{M}_{\epsilon_i}$} (3.05, -1.95);

% =====================================================================
%  USER AGENT (center): sanitizer + sanitized cache + interaction history
% =====================================================================

% Sanitizer
\node[sanitizer] (san) at (4.0, -0.4) {Sanitizer\\\scriptsize $\mathcal{M}_{\epsilon_i}$};

% Sanitized cache (3 noised roots)
\node[cached] (xh1) at (5.5, -1.5) {$\hat{x}_1$};
\node[cached] (xh2) at (6.4, -1.5) {$\hat{x}_2$};
\node[cached] (xh3) at (7.3, -1.5) {$\hat{x}_3$};

\begin{scope}[on background layer]
    \node[draw=blue!50!black, thick, fill=blue!4, rounded corners=3pt,
          fit=(xh1)(xh3), inner xsep=5pt, inner ysep=4pt] (cache) {};
\end{scope}
\node[smalllabel, text=blue!50!black, anchor=north] at (cache.south) {sanitized cache};

% Sanitizer -> cache
\draw[arr, blue!60!black] (san.south) -- ([yshift=2pt]cache.north);

% Interaction history (4 turns, mirrors original fig:dag_overview)
\node[derived] (q1) at (4.55, -2.85) {$\hat{x}_1$};
\node[derived] (q2) at (5.45, -2.85) {$\hat{x}_1$};
\node[derived] (q3) at (6.7, -2.85) {$f(\hat{x}_2, \hat{x}_3)$};
\node[tgt]     (qT) at (7.85, -2.85) {$T(\hat{X})$};

\node[smalllabel, anchor=north] at (q1.south) {turn 1};
\node[smalllabel, anchor=north] at (q2.south) {turn 2};
\node[smalllabel, anchor=north] at (q3.south) {turn 3};
\node[smalllabel, anchor=north] at (qT.south) {turn 4};

% Same-value annotation across turns 1 and 2 (matches original)
\draw[same] ([xshift=-3pt,yshift=2pt]q1.north west) rectangle ([xshift=3pt,yshift=-12pt]q2.south east);

% Cache -> interactions
\draw[arr, black!30] (xh1.south) -- (q1.north);
\draw[arr, black!30] (xh1.south) -- (q2.north);
\draw[arr, black!30] (xh2.south) -- (q3.north);
\draw[arr, black!30] (xh3.south) -- (q3.north);
\draw[arr, black!30] (xh1.south) -- (qT.north);
\draw[arr, black!30] (xh2.south) -- (qT.north);
\draw[arr, black!30] (xh3.south) -- (qT.north);

% Agent enclosure
\begin{scope}[on background layer]
    \node[draw=black!40, fill=black!2, rounded corners=4pt, dotted, thick,
          fit=(san)(cache)(q1)(qT), inner xsep=10pt, inner ysep=14pt] (agentbox) {};
\end{scope}
\node[font=\small\bfseries, text=black!60, anchor=south west]
    at ([xshift=2pt, yshift=2pt]agentbox.north west) {\sysname{} (user agent)};

% =====================================================================
%  TRUST BOUNDARY
% =====================================================================

\draw[boundary] (10.0, 0.7) -- (10.0, -3.6);
\node[font=\scriptsize\bfseries, text=red!50!black, anchor=south, rotate=90]
    at (10.0, -3.6) {trust boundary};

% =====================================================================
%  ADVERSARIAL SERVICE (right)
% =====================================================================

\node[service] (svc) at (12.0, -1.95) {Adversarial\\Service $\mathcal{A}$};
\node[smalllabel, text=red!50!black, anchor=north, align=center]
    at (svc.south) {logs / colludes /\\breached};

% =====================================================================
%  CROSS-BOUNDARY ARROWS (both straight)
%  Top: queries / functions  (service -> agent)
%  Bottom: interactions      (agent  -> service)
% =====================================================================

\draw[flow, red!50!black, dashed]
    ([yshift=8pt]svc.west) --
    node[smalllabel, above, text=red!50!black] {queries / functions $f$}
    ([yshift=8pt]agentbox.east |- svc.west);

\draw[flow, blue!50!black]
    ([yshift=-8pt]agentbox.east |- svc.west) --
    node[smalllabel, below, text=blue!50!black] {interactions (cached / post-processed)}
    ([yshift=-8pt]svc.west);

\end{tikzpicture}

\caption{Overview of \sysname{} in an agentic deployment. At initialization, the user passes its private values $X = (x_1, x_2, x_3)$ and a chosen mechanism $\mathcal{M}_{\epsilon_i}$ to the user agent. \sysname{} sanitizes each root once, stores the result in the sanitized cache, and discards the originals. Subsequent interactions with the adversarial service are answered entirely from the cache: repeated requests for the same root return the identical value (turns 1--2), derived computations are deterministic post-processing of cached roots (turn 3), and the target incurs no additional privacy cost (turn 4). The trust boundary separates the user agent (which holds noised values) from the adversarial service, which may log, collude with other services, or be breached. By the post-processing theorem, the adversary's view is no more informative than the initial root sanitization, regardless of which functions it requests.}
\label{fig:rootguard_overview}
\end{figure*}